
\documentclass[UTF8,11pt,a4paper]{ctexart}
\usepackage[a4paper,left=23mm,right=23mm,top=24mm,bottom=24mm]{geometry}
\usepackage{fontspec}
\setmainfont{DejaVu Sans}
\setCJKmainfont{Noto Serif CJK SC}
\usepackage{amsmath,amssymb}
\usepackage{xcolor}
\usepackage[colorlinks=true,linkcolor=blue!45!black,unicode=true]{hyperref}
\usepackage{enumitem}
\usepackage{needspace}
\usepackage{fancyhdr}
\setlength{\parindent}{2em}
\setlength{\parskip}{0.25em}
\setlist[enumerate]{leftmargin=2.6em,itemsep=0.4em,topsep=0.5em}
\pagestyle{fancy}
\fancyhf{}
\fancyhead[L]{\small 三千弱水 · 大学化学高阶综合题库}
\fancyhead[R]{\small 待专家审校稿}
\fancyfoot[C]{\thepage}
\setlength{\headheight}{15pt}
\sloppy

\begin{document}
\begin{center}
{\LARGE\bfseries 大学化学高阶综合题 · 第2辑 (12题)}\\[0.8em]
{\large 题目、分步解析与数值核验}\\[0.5em]
{\small \textcolor{red!65!black}{审校稿：未经化学专家核准；禁止直接视为正式试卷或出版物}}
\end{center}
\vspace{1em}
\noindent 本册从仓库结构化 JSONL 自动生成。增补和修改应只编辑源题目数据，并重新执行统一编译。
\par 机械方程残差检查只能说明指定模型的内部数值一致，不保证模型现实可行性。第六辑另须完成跨辑模型去重审查。
\tableofcontents
\clearpage
\section*{第2辑 · 题目及解答}
\addcontentsline{toc}{section}{第2辑}
\subsection*{无机化学}
\addcontentsline{toc}{subsection}{第2辑 / 无机化学}
\Needspace{7\baselineskip}
\subsubsection*{UC-ADV2-I01 · 竞争配体与选择性硫化沉淀的分段平衡反演}
\noindent\textbf{题目}\par
假想二价金属M、N在1.000 L缓冲溶液中分析浓度分别为0.01000、0.02000 mol/L，单负配体L总浓度0.02500 mol/L。仅形成ML+、NL+，条件稳定常数βM=400、βN=120 L/mol，忽略质子化与其它络合。初态无固相。(1)求初始游离L、M2+、N2+。(2)固定游离[S2-]=1.00×10\textasciicircum{}-5 mol/L并允许仅MS(s)沉淀，Ksp(MS)=1.00×10\textasciicircum{}-8、Ksp(NS)=1.00×10\textasciicircum{}-6；重新求游离L、两金属的游离/络合浓度与沉淀M量。(3)验证NS不沉淀、元素及配体守恒，并说明沉淀后为什么不能沿用初始L。\par
\medskip\noindent\textbf{分步解答}
\begin{enumerate}
\item 初态配体守恒：CL=l+CM βM l/(1+βM l)+CN βN l/(1+βN l)；二价游离浓度m0=CM/(1+βM l), n0=CN/(1+βN l)。
\item 左侧是l的严格递增函数，在0与CL分别跨过CL，故存在唯一物理解；二分法得l0=0.0077775022 mol/L，m0=0.00243249766 mol/L，n0=0.0103450045 mol/L。
\item 加入指定的自由S2−后，若MS(s)存在，游离m=Ksp(MS)/[S2−]=0.001 mol/L；它不能继续等于m0。
\item 重新建立配体守恒：l+βM m l+CN βN l/(1+βN l)=CL。该函数严格递增，解得l=0.0100479839 mol/L。
\item 于是[ML+]=βM m l=0.00401919358 mol/L，[N2+]=CN/(1+βN l)=0.00906717751 mol/L，[NL+]=CN−[N2+]=0.0109328225 mol/L。
\item M沉淀量=CM−[M2+]−[ML+]=0.00498080642 mol/L；N的离子积=[N2+][S2−]=9.06717751e-8<0.000001，故NS不沉淀。
\item 回代：游离L+络合所占L=0.025 mol/L；M、N各自总量均等于投料量；若沉淀量算成负值，则事先的相假设必须放弃。
\end{enumerate}
\noindent\textbf{答案摘要}\par
答案1 = 0.0077775022 mol/L；答案2 = 0.00243249766 mol/L；答案3 = 0.0103450045 mol/L；答案4 = 0.0100479839 mol/L；答案5 = 0.001 mol/L；答案6 = 0.00401919358 mol/L；答案7 = 0.00906717751 mol/L；答案8 = 0.0109328225 mol/L；答案9 = 0.00498080642 mol/L\par
\noindent\textbf{模型前提：}
模型为假想离子体系；固定S2−和pH由外部缓冲储液维持，电中性由未逐项列出的支持电解质反离子满足；所有常数均以题定浓度近似。\par
\medskip\hrule\medskip
\Needspace{7\baselineskip}
\subsubsection*{UC-ADV2-I02 · 混合价钴钙钛矿的氧空位、晶胞密度与闭孔率}
\noindent\textbf{题目}\par
某理想化SrCoO\_(3−δ)钙钛矿每立方晶胞含1个化学式单位，X射线测得a=0.3890 nm，独立价态谱给出Co4+摩尔分数0.700（其余Co3+），晶体宏观体密度5.180 g/cm³。假设Sr2+、O2−，唯一点缺陷是氧空位，晶胞体积不受缺陷影响；Sr、Co、O摩尔质量分别为87.62、58.933、15.999 g/mol。(1)由电中性求δ。(2)求无孔含氧空位晶体的理论密度。(3)解释理论与实测密度差，并估算闭孔体积分数。(4)说明把氧空位与孔洞混淆会引起的偏差。\par
\medskip\noindent\textbf{分步解答}
\begin{enumerate}
\item 设Co4+占比x=0.700，平均Co氧化态=3+x=3.7。
\item 电荷守恒：+2+(3+x)−2(3−δ)=0，即δ=(1−x)/2=0.15，故化学式SrCoO\_2.85。
\item 每化学式单位摩尔质量M=87.62+58.933+(3−δ)×15.999=192.15015 g/mol。
\item 立方晶胞体积V=(0.3890×10\textasciicircum{}-7 cm)\textasciicircum{}3=5.8863869e-23 cm³；质量m=M/NA。
\item 无孔理论密度ρ0=M/(NA a³)=5.42052088 g/cm³。
\item 如仅有不可进入的闭孔且骨架密度ρ0，ρbulk=(1−φ)ρ0，故φ=1−5.180/5.42052088=4.43722813\%。
\item 正向核对：含氧数×(−2)恰抵Sr与Co总正电；反向核对：(1−φ)ρ0=5.180 g/cm³。
\end{enumerate}
\noindent\textbf{答案摘要}\par
答案1 = 0.15 per formula unit；答案2 = 5.42052088 g/cm³；答案3 = 0.0443722813 fraction\par
\noindent\textbf{模型前提：}
Co仅有+3/+4价、只有氧位空位、缺陷不改晶胞常数，宏观密度下降剩余部分仅来自闭孔；不是实际样品的未经验证结论。\par
\medskip\hrule\medskip
\Needspace{7\baselineskip}
\subsubsection*{UC-ADV2-I03 · 自旋交叉磁化率反演焓熵与半转换温度}
\noindent\textbf{题目}\par
理想独立Fe(II)自旋交叉中心，仅有低自旋S=0和高自旋S=2两态，忽略轨道贡献及相互作用，高自旋spin-only有效磁矩μH=√24 μB，低自旋磁矩取0。实验按Curie权重得到280.0 K的μeff=√6 μB、320.0 K的μeff=√14.4 μB。设高/低态平衡比K=x/(1−x)=exp(ΔS/R−ΔH/(RT))，ΔH、ΔS与温度无关。(1)求高自旋摩尔分数。(2)反演ΔH、ΔS。(3)求50\%转换温度及其存在条件。(4)检验升温方向是否与焓符号一致。\par
\medskip\noindent\textbf{分步解答}
\begin{enumerate}
\item 磁化率线性加和的是平方磁矩：μeff²=x μH²+(1−x)μL²=24x μB²。
\item 280 K时x1=6/24=0.25；320 K时x2=14.4/24=0.6。
\item K1=x1/(1−x1)=0.333333333，K2=x2/(1−x2)=1.5。
\item 两温度van't Hoff式相减：ΔH=R ln(K2/K1)/(1/T1−1/T2)=28.0125334 kJ/mol。
\item 再由lnK1=ΔS/R−ΔH/(RT1)，得到ΔS=RlnK1+ΔH/T1=90.9103915 J/(mol K)。
\item 半转换时K=1，0=ΔS−ΔH/T50，故T50=ΔH/ΔS=308.13346 K（要求ΔS、ΔH同号且T50>0）。
\item 正向回代μeff1=2.44948974 μB、μeff2=3.79473319 μB，且升温高自旋分数增大，与ΔH>0的吸热高自旋态一致。
\end{enumerate}
\noindent\textbf{答案摘要}\par
答案1 = 0.25 fraction；答案2 = 0.6 fraction；答案3 = 28012.5334 J/mol；答案4 = 90.9103915 J/(mol K)；答案5 = 308.13346 K\par
\noindent\textbf{模型前提：}
有效磁矩按理想Curie权重的平方平均；不考虑零场分裂、交换耦合、磁各向异性和热容差。\par
\medskip\hrule\medskip
\subsection*{分析化学}
\addcontentsline{toc}{subsection}{第2辑 / 分析化学}
\Needspace{7\baselineskip}
\subsubsection*{UC-ADV2-A01 · 未知基质灵敏度与标准加入的双波长联合反演}
\noindent\textbf{题目}\par
未知样含X、Y两种发色物，1.00 cm池，两波长吸光系数矩阵[[8000,2000],[3000,7000]] L mol\textasciicircum{}-1 cm\textasciicircum{}-1。由于基质效应，两种物的整体信号均乘未知公共抑制因子q∈(0,1)，测得原样(A1,A2)=(0.2952,0.2132)。向等终体积的平行样品标准加入ΔcX=1.000×10\textasciicircum{}-5 mol/L的X，保持基质因子不变，测得(0.3608,0.2378)。(1)同时求q、cX、cY。(2)检验加标两波长是否自洽。(3)说明单凭未加标的两个吸光度为什么无法确定三个未知量，并写出模型失效的判断方式。\par
\medskip\noindent\textbf{分步解答}
\begin{enumerate}
\item 未加标向量A=q ε(cX,cY)\textasciicircum{}T，加标后A'=A+q ΔcX (8000,3000)\textasciicircum{}T。
\item 第一波长给q=(A1'−A1)/(8000ΔcX)=0.82；第二波长独立求q=(A2'−A2)/(3000ΔcX)=0.82，两者吻合。
\item 除去公共基质因子后b=A/q=(0.36,0.26)。
\item 线性方程ε c=b，行列式=8000×7000−2000×3000=50000000，不为零。
\item Cramer's rule：cX=(7000b1−2000b2)/det=0.00004 mol/L，cY=(8000b2−3000b1)/det=0.00002 mol/L。
\item 回代预测未加标(0.2952,0.2132)及加标(0.3608,0.2378)。
\item 只有两条未加标方程而未知量为q,cX,cY三者，系统欠定；加标增量独立识别q，若两波长增量求得q不同，公共抑制假设应被否定。
\end{enumerate}
\noindent\textbf{答案摘要}\par
答案1 = 0.82 dimensionless；答案2 = 4e-05 mol/L；答案3 = 2e-05 mol/L\par
\noindent\textbf{模型前提：}
两物在两波长的基质抑制是同一个正比例常数；加标前后终体积相同、无反应、无非线性吸收。\par
\medskip\hrule\medskip
\Needspace{7\baselineskip}
\subsubsection*{UC-ADV2-A02 · 封闭碳酸体系加酸后的精确电荷平衡与物种分布}
\noindent\textbf{题目}\par
一个100.0 mL封闭水样起初含NaHCO3 0.0200 mol/L和Na2CO3 0.0100 mol/L，加入0.00150 mol HCl后定容120.0 mL，封闭体系不逸散CO2。25°C时取Ka1=10\textasciicircum{}-6.35、Ka2=10\textasciicircum{}-10.33、Kw=10\textasciicircum{}-14，理想浓度活度近似。求最终无机碳总浓度、[Na+]、[Cl−]、准确的pH以及H2CO3*、HCO3−、CO3\textasciicircum{}2−各物种浓度。不得直接把Henderson–Hasselbalch近似当作精确电荷平衡。\par
\medskip\noindent\textbf{分步解答}
\begin{enumerate}
\item 无机碳守恒CT=[(0.0200+0.0100)×0.1000]/0.1200=0.025 mol/L。
\item 钠与氯分别守恒：[Na+]=(0.0200+2×0.0100)×0.1000/0.1200=0.0333333333 M，[Cl−]=0.00150/0.1200=0.0125 M。
\item 令h=[H+]，D=h²+Ka1 h+Ka1 Ka2，则[H2CO3*]=CT h²/D，[HCO3−]=CT Ka1 h/D，[CO3²−]=CT Ka1 Ka2/D。
\item 电荷守恒式：[Na+]+h=[Cl−]+[HCO3−]+2[CO3²−]+Kw/h；等价求根F(pH)= [HCO3−]+2[CO3²−]+Kw/h−h−([Na+]−[Cl−])=0。
\item 在pH为4到11区间以单调物种响应数值二分，得pH=7.0473843。
\item 代回[H2CO3*]=0.00417754503 M、[HCO3−]=0.0208115985 M、[CO3²−]=0.0000108564975 M。
\item 检查碳守恒：三形态和=0.025 M，电荷守恒右侧净负电荷=0.0208333333 M，等于Na−Cl=0.0208333333 M。
\end{enumerate}
\noindent\textbf{答案摘要}\par
答案1 = 0.025 mol/L；答案2 = 0.0333333333 mol/L；答案3 = 0.0125 mol/L；答案4 = 7.0473843 pH；答案5 = 0.00417754503 mol/L；答案6 = 0.0208115985 mol/L；答案7 = 1.08564975e-05 mol/L\par
\noindent\textbf{模型前提：}
封闭、定容后处于平衡；H2CO3*包含水合CO2的合并定义；给定常数按浓度条件使用，忽略活度系数与空气CO2交换。\par
\medskip\hrule\medskip
\Needspace{7\baselineskip}
\subsubsection*{UC-ADV2-A03 · 共沉积电量—质量双守恒反演金属组成}
\noindent\textbf{题目}\par
100.0 mL水样只含Cu2+与Ag+两种可沉积金属离子，恒电流0.1500 A电沉积2656.84248 s，电流效率92.0\%，电极净增质量0.18161356 g（无夹杂），F=96485.33212 C/mol。反应Cu2++2e→Cu、Ag++e→Ag；MCu=63.546、MAg=107.8682 g/mol。(1)联立法拉第电荷衡算与沉积质量求两金属的初始摩尔数及浓度。(2)计算理论用于沉积的电量、损失电量。(3)讨论若两种离子都错按二电子还原为何不能可靠反演。\par
\medskip\noindent\textbf{分步解答}
\begin{enumerate}
\item 总通过电量Q=It=0.1500×2656.84248=398.526372 C；有效沉积电量Q沉=0.920Q=366.644262 C。
\item 设nCu=u、nAg=v，法拉第定律(2u+v)=Q沉/F=0.0038 mol电子；称重方程63.546u+107.8682v=0.18161356 g。
\item 消去v：u=[(Q沉/F)−m/MAg]/[2−MCu/MAg]=0.0015 mol。
\item v=(Q沉/F)−2u=0.0008 mol。
\item 以样品体积0.1000 L折算：[Cu2+]=0.015 mol/L，[Ag+]=0.008 mol/L。
\item 副反应电量Q损=(1−0.920)Q=31.8821097 C，不能计入金属沉积量。
\item 逆向验算：理论m=0.18161356 g，理论沉积电荷=366.644262 C，两者恢复观测。错误地令Ag+也消耗2e会改变电量矩阵的秩与正确组成。
\end{enumerate}
\noindent\textbf{答案摘要}\par
答案1 = 0.0015 mol；答案2 = 0.0008 mol；答案3 = 0.015 mol/L；答案4 = 0.008 mol/L；答案5 = 366.644262 C；答案6 = 31.8821097 C\par
\noindent\textbf{模型前提：}
金属定量回收到同一电极且无合金化引起的重量差；电流效率已独立校准，电解时无化学溶解或其它阳离子沉积。\par
\medskip\hrule\medskip
\subsection*{有机化学}
\addcontentsline{toc}{subsection}{第2辑 / 有机化学}
\Needspace{7\baselineskip}
\subsubsection*{UC-ADV2-O01 · 平行—串联反应网络的中间产物选择性最优淬灭}
\noindent\textbf{题目}\par
恒容恒温下，假想有机反应网络A→B（k1=0.110 h\textasciicircum{}-1）、A→D（副反应k3=0.030 h\textasciicircum{}-1）、B→C（k2=0.045 h\textasciicircum{}-1），均为不可逆一级反应，初始仅A，A0=1.000 mol。各路径均为1:1计量，D和C稳定。(1)写四物种浓度/摩尔分数随时间的显式解。(2)求B最大分离时刻及四个物种的分数。(3)给出最优时A→B的总反应选择性上限，以及质量守恒证明。\par
\medskip\noindent\textbf{分步解答}
\begin{enumerate}
\item 记K=k1+k3=0.140 h\textasciicircum{}-1，微分方程a'=−Ka，b'=k1a−k2b，d'=k3a，c'=k2b，初值(a,b,c,d)=(1,0,0,0)。
\item a(t)=exp(−Kt)；b(t)=k1/(K−k2)[exp(−k2t)−exp(−Kt)]。
\item d(t)=(k3/K)[1−exp(−Kt)]；c(t)=1−a−b−d（直接由总物质守恒得到）。
\item 令b'(t)=0：k2 exp(−k2t)=K exp(−Kt)，所以t*=ln(K/k2)/(K−k2)=11.9471572 h。
\item 该时刻(a,b,c,d)=(0.187757882,0.458963712,0.179226523,0.174051882)；因此最大可瞬时分离B为0.458963712 mol。
\item 如果B不继续消耗，A一旦反应进入B通道的最终分支概率为k1/(k1+k3)=0.785714286，这是对B分离总收率的绝对上限；本题在t*之前存在B→C损失，所以实际b*严格小于该上限。
\item 正向验证a+b+c+d=1；并且b'(t*)=k1a−k2b=3.46944695e-18 mol h\textasciicircum{}-1，且b(0)=b(∞)=0，故峰值唯一。
\end{enumerate}
\noindent\textbf{答案摘要}\par
答案1 = 11.9471572 h；答案2 = 0.187757882 fraction；答案3 = 0.458963712 fraction；答案4 = 0.179226523 fraction；答案5 = 0.174051882 fraction\par
\noindent\textbf{模型前提：}
假想单向一级动力学、均一混合，忽略可逆反应、分子量改变造成的质量不等价；分数按1:1分子摩尔计。\par
\medskip\hrule\medskip
\Needspace{7\baselineskip}
\subsubsection*{UC-ADV2-O02 · 有限物料酯化反应中间脱水与二次平衡转化}
\noindent\textbf{题目}\par
理想恒温液相酯化A+B⇌E+H2O，按活度≈摩尔浓度/标准浓度处理，K=4.000，体积近似恒定。最初nA=1.000、nB=1.200、nE=0、nW=0.0500 mol。(1)求第一次平衡转化ξ1及全部组分。(2)随后仅从体系除去0.1500 mol水，反应物、酯不损失，重新封闭达到第二次平衡，求追加反应进度ξ2及总转化率。(3)两次完整物料守恒反算，解释为何必须把第一次形成的E、W带入第二次平衡。\par
\medskip\noindent\textbf{分步解答}
\begin{enumerate}
\item 设第一次正向反应进度ξ1，反应商Q=(nE nW)/(nA nB)（体积的同次幂消去）。
\item 平衡方程ξ1(0.0500+ξ1)=4(1−ξ1)(1.2−ξ1)，物理根须满足0<ξ1<1；求得ξ1=0.716300455 mol。
\item 第一次平衡组成：(nA,nB,nE,nW)=(0.283699545,0.483699545,0.716300455,0.766300455) mol。
\item 仅移除0.1500 mol水后、二次反应前nW=0.616300455 mol，其余组成不变。
\item 令后续追加转化ξ2，第二次平衡满足(E+ξ2)(W+ξ2)=4(A−ξ2)(B−ξ2)，并限定0≤ξ2≤min(A,B)。
\item 二次物理根ξ2=0.0248272013 mol，最终总酯形成量ξ1+ξ2=0.741127656 mol，相对最初A的总酯化转化率=74.1127656\%。
\item 检查初始A：nA最终+ξ1+ξ2=1 mol；最终水量nW=0.641127656 mol，与移水操作及反应生成守恒一致。
\end{enumerate}
\noindent\textbf{答案摘要}\par
答案1 = 0.716300455 mol；答案2 = 0.0248272013 mol；答案3 = 0.741127656 mol；答案4 = 0.258872344 mol；答案5 = 0.641127656 mol\par
\noindent\textbf{模型前提：}
完整计量反应，恒温、溶液理想及体积对反应商抵消，移除水时未带走其它组分；无副反应。\par
\medskip\hrule\medskip
\Needspace{7\baselineskip}
\subsubsection*{UC-ADV2-O03 · 反应活性比控制下的瞬时共聚组成与进料漂移}
\noindent\textbf{题目}\par
两单体M1、M2按终端自由基共聚模型，反应活性比分别r1=0.200、r2=2.000，初始液相M1的摩尔分数f1,0=0.400；不补料，忽略体积变化和链转移对瞬时组成的影响。Mayo–Lewis瞬时共聚M1摩尔分数F1=(r1 f1²+f1(1−f1))/(r1 f1²+2f1(1−f1)+r2(1−f1)²)。(1)求初始F1。(2)推导随总单体转化率X变化的进料f1微分方程。(3)数值积分至X=0.400，求剩余进料f1、此时瞬时F1及累计聚合物平均M1组分。(4)做M1和总单体两种物料守恒核验。\par
\medskip\noindent\textbf{分步解答}
\begin{enumerate}
\item 由给定Mayo–Lewis式，初始f1=0.400时F1,0=0.220779221，与初始进料不相同，产生组成漂移。
\item 初始总单体N0，转化率X时剩余单体N=N0(1−X)，其中M1数量为N0(1−X)f1。
\item 新消耗d(N0X)个单体时M1消耗F1 d(N0X)，所以d[(1−X)f1]/dX=−F1(f1)。
\item 展开得df1/dX=[f1−F1(f1)]/(1−X)，初值f1(0)=0.400；这是一阶非线性微分方程。
\item 用四阶Runge–Kutta步长ΔX=0.0004积分至X=0.400，得f1末=0.500530342，此刻瞬时F1末=0.286075161。
\item M1总守恒：初始0.4N0，末剩0.6f1末N0，因此聚合物中累计M1= [0.4−0.6f1末]N0；其累计平均摩尔分数Fbar=[0.4−0.6f1末]/0.4=0.249204487。
\item 将积分步长减半到0.0002后f1末=0.500530342，与原解一致到数值精度；Fbar与单体总消耗0.400N0相容。
\end{enumerate}
\noindent\textbf{答案摘要}\par
答案1 = 0.220779221 fraction；答案2 = 0.500530342 fraction；答案3 = 0.286075161 fraction；答案4 = 0.249204487 fraction\par
\noindent\textbf{模型前提：}
终端共聚反应活性比定值、两种单体无其它损耗或补料；题目只描述组成漂移，不声称给出聚合物分子量分布。\par
\medskip\hrule\medskip
\subsection*{物理化学}
\addcontentsline{toc}{subsection}{第2辑 / 物理化学}
\Needspace{7\baselineskip}
\subsubsection*{UC-ADV2-P01 · 绝热反应釜中热容差、温度与平衡转化的耦合非线性求解}
\noindent\textbf{题目}\par
绝热恒压反应釜中初始有1.000 mol纯A以及不参与反应的热容为100.0 J/K的惰性材料，初温298.15 K。反应A⇌B，ΔrH°298=−40.00 kJ/mol，298.15 K平衡常数K=2.400；A、B定压摩尔热容分别为100和120 J/(mol K)，常数。理想等摩尔数气相单反应（K=nB/nA），忽略流动和机械外功之外的额外热源。(1)建立绝热焓守恒得到T(ξ)。(2)用Kirchhoff与van't Hoff积分写K(T)。(3)联立求最终T、ξ并独立检查焓守恒与反应商。\par
\medskip\noindent\textbf{分步解答}
\begin{enumerate}
\item 反应焓ΔH(T)=H0+ΔCp(T−T0)，ΔCp=120−100=20 J/(mol K)。
\item 以初始A+惰性材料为参考，绝热恒压焓平衡：(200+20ξ)(T−298.15)−40000ξ=0。
\item 于是T(ξ)=298.15+40000ξ/(200+20ξ) K。
\item 积分dlnK/dT=ΔH(T)/(RT²)，得ln[K(T)/K0]=H0/R(1/T0−1/T)+(ΔCp/R)[ln(T/T0)+T0/T−1]。
\item 理想气体等摩尔数体系K(T)=ξ/(1−ξ)，将T(ξ)代入，构成0<ξ<1上的单变量非线性求根问题。
\item 经二分法得到ξ=0.229384325，T=342.998119 K，K(T)=0.297663715，并由ξ/(1−ξ)=0.297663715 验证平衡。
\item 绝热残差(200+20ξ)(T−298.15)−40000ξ=-5.45696821e-12 J；升温降低放热反应平衡常数，负反馈使ξ显著低于恒温时K0/(1+K0)。
\end{enumerate}
\noindent\textbf{答案摘要}\par
答案1 = 0.229384325 mol reaction per initial mol A；答案2 = 342.998119 K；答案3 = 0.297663715 dimensionless\par
\noindent\textbf{模型前提：}
焓用恒压热容，容器近绝热，理想等计量系数气相使K仅为摩尔数比，ΔCp常数、无相变。\par
\medskip\hrule\medskip
\Needspace{7\baselineskip}
\subsubsection*{UC-ADV2-P02 · 有机相二聚诱导的非线性两级液液萃取}
\noindent\textbf{题目}\par
1.000 L水相含单体型溶质A共0.01000 mol，依次与两份新鲜0.2000 L有机相平衡萃取。水相只含单体A，单体分配比[A]有机/[A]水=P=3.000；有机相中2A⇌A2二聚，形成常数Kdim=[A2]有机/[A]有机²=40.0 L/mol。忽略体积改变、乳化和非理想性。(1)建立含二聚体的准确物料平衡并求第一级后水相浓度。(2)第二级使用新鲜有机相，求最终浓度、总萃取率。(3)分别求两级提走的A结构单元物质的量，并检验全流程守恒；不可将有机相二聚体当作只含一个A单元。\par
\medskip\noindent\textbf{分步解答}
\begin{enumerate}
\item 设一次萃取后水相单体浓度c；有机相单体浓度Pc、二聚体浓度Kdim(Pc)²。
\item 以A结构单元计，单次物料守恒为n输入=V水 c+V有[P c+2Kdim(P c)²]，即n=c+0.2000(3c+720c²)。
\item 第一级由n0=0.01000 mol求唯一非负根c1=0.00445986465 mol/L；第一级有机相提走n有,1=0.01000−c1=0.00554013535 mol。
\item 第二级有机相全新，初始水相只剩n输入=c1×1.000 L，因此必须重新求c2+0.2000(3c2+720c2²)=c1。
\item 得到c2=0.00230799834 mol/L，第二级提走n有,2=c1−c2=0.00215186632 mol。
\item 总萃取率η=(n0−1.000c2)/n0=76.9200166\%；不能用常数分配比的几何级数，因为二聚使萃取量非线性依赖浓度。
\item 回代全流程：n有,1+n有,2+n水,末=0.01 mol，与最初0.01000 mol相同；两个二次式在c≥0区间严格递增，故各有唯一物理解。
\end{enumerate}
\noindent\textbf{答案摘要}\par
答案1 = 0.00445986465 mol/L；答案2 = 0.00230799834 mol/L；答案3 = 0.00554013535 mol A；答案4 = 0.00215186632 mol A；答案5 = 0.769200166 fraction\par
\noindent\textbf{模型前提：}
只允许有机相A2二聚，水相无二聚或其他反应，P与Kdim为浓度基准常数；萃取每级达平衡。\par
\medskip\hrule\medskip
\Needspace{7\baselineskip}
\subsubsection*{UC-ADV2-P03 · 有限NaCl储液池混合的最大可逆电功与自由能审计}
\noindent\textbf{题目}\par
298.15 K下，两份理想完全解离的NaCl水溶液分别为储液池1：V1=1.000 L、c1=0.02000 mol/L；储液池2：V2=2.000 L、c2=0.005000 mol/L。两池通过理想可逆盐浓差电能装置充分交换，最终均匀并保持总体积不变。两种离子遵循理想溶液活度模型、无溶剂交换功和液接损耗。(1)求最终浓度。(2)从各池初末盐离子数推导总Gibbs混合自由能差。(3)求最大理论可回收电功，并说明为何不能把初始电压当作全程定值直接乘一个固定电量。(4)核验正负号、总离子数量和零梯度极限。\par
\medskip\noindent\textbf{分步解答}
\begin{enumerate}
\item NaCl分析量守恒：n总=V1c1+V2c2=0.03 mol，最终c\_f=n总/(V1+V2)=0.01 mol/L。
\item NaCl每个化学式单位分解为两个单价离子，理想化学势μ\_i=μ\_i°+RTln(c/c°)，标准态项在守恒前后相消。
\item 恒温下总自由能变化ΔG=2RT[V1c1 ln(cf/c1)+V2c2 ln(cf/c2)]，其中V以L、c以mol/L得到mol。
\item 代入298.15 K，ΔG=-34.3656415 J；可逆最大输出电功Wmax=−ΔG=34.3656415 J。
\item 浓差耗散会令ΔG<0，而纯混合不回收功；若以可逆装置回收，理论可输出正功34.3656415 J，实际更低。
\item 离子守恒：总Na+与Cl−各0.03 mol；若c1=c2，则每个ln(cf/ci)=0，最大功为零。
\item 浓差电池电势E随两池浓度和转运进度变化，故不能使用初态E×总转运电量代替积分−ΔG。
\end{enumerate}
\noindent\textbf{答案摘要}\par
答案1 = 0.01 mol/L；答案2 = -34.3656415 J；答案3 = 34.3656415 J\par
\noindent\textbf{模型前提：}
仅计算固定温度、溶液体积不变且无其他机械/混合溶剂功的盐离子理想自由能差；不把真实膜电位与不可逆损耗忽略的结果冒充实验值。\par
\medskip\hrule\medskip
\end{document}
